% =====================================================================
\section{Sources}
{\small
\textbf{Verified while preparing the brochure} (tag \tP):
\begin{itemize}
\item Saha~S. et al. A systematic review of the prevalence of schizophrenia. PLoS Medicine, 2005.
\item Hilker~R. et al. Heritability of schizophrenia and schizophrenia spectrum based on the nationwide Danish twin register. Biological Psychiatry, 2018 (as reported in the press release).
\item Krystal~J.\,H. et al. Subanesthetic effects of the noncompetitive NMDA antagonist, ketamine, in humans. Archives of General Psychiatry, 1994.
\item Kapur~S. et al. Relationship between dopamine D2 occupancy, clinical response, and side effects. American Journal of Psychiatry, 2000.
\item Howes~O.\,D. et al. The nature of dopamine dysfunction in schizophrenia and what this means for treatment. Archives of General Psychiatry, 2012.
\item A study of glutathione turnover in the mouse brain, PubMed record 9165485.
\item Harkness~J.\,H. et al. Diurnal changes in perineuronal nets and parvalbumin neurons in the rat medial prefrontal cortex. bioRxiv preprint, 2020; published 2021.
\item Molecular and cellular rhythms in excitatory and inhibitory neurons in the mouse prefrontal cortex. bioRxiv preprint, 2024.
\item Schizophrenia and catatonia: from ICD-10 to ICD-11. Review article, Springer, 2025 (authors not verified). The list of ICD-11 codes of group 6A2 is from the code reference.
\item Hashimoto~T. et al. Relationship of brain-derived neurotrophic factor and its receptor TrkB to altered inhibitory prefrontal circuitry in schizophrenia. Journal of Neuroscience, 2005.
\item Fernandes~B.\,S. et al. Peripheral brain-derived neurotrophic factor in schizophrenia and the role of antipsychotics: meta-analysis and implications. Molecular Psychiatry, 2015.
\item Durany~N., Thome~J. Neurotrophic factors and the pathophysiology of schizophrenic psychoses. European Psychiatry, 2004; Angelucci~F. et al. BDNF in schizophrenia, depression and corresponding animal models. Molecular Psychiatry, 2005.
\item Jaggar~M., Vaidya~V.\,A. 5-HT2A receptors and BDNF regulation: implications for psychopathology. Book chapter, Springer, 2018 (from the abstract).
\item Papadia~S. et al. Synaptic NMDA receptor activity boosts intrinsic antioxidant defenses. Nature Neuroscience, 2008 (from a secondary account).
\end{itemize}
\textbf{Read by the authors of the project} (tag \tS):
\begin{itemize}
\item Hardingham~G.\,E., Do~K.\,Q. Linking early-life NMDAR hypofunction and oxidative stress in schizophrenia pathogenesis. Nature Reviews Neuroscience, 2016.
\item Cuenod~M. et al. Review of the vicious circles through oxidative stress, 2021, and the reply to criticism (Molecular Psychiatry, 2022).
\item Selvaggio~G., Coelho~P., Salvador~A. A model of the peroxiredoxin--thioredoxin system. Redox Biology, 2018.
\item Volman~V., Behrens~M.\,M., Sejnowski~T.\,J. A network model with loss of parvalbumin and GAD67, 2011.
\item Kinney~J.\,W. et al., 2006 --- ketamine in culture; Behrens~M.\,M. et al., 2008 --- ketamine in mice; Baxter~P.\,S. et al., 2015 --- the receptor and glutathione synthesis; Brennan~A.\,M. et al., 2009 --- the receptor and Nox2.
\item Kil~I.\,S. et al. Circadian oscillation of sulfiredoxin in the mitochondria. Molecular Cell, 2015; del~Olmo~M. et al. A robust model for circadian redox oscillations. International Journal of Molecular Sciences, 2019.
\item Edgar~R.\,S. et al., 2012 --- daily cycles of the peroxiredoxins; Wang~T.\,A. et al., 2012 --- oxidative state and excitability in the master clock; Kinoshita~C. et al., 2014 --- the daily rhythm of glutathione.
\item Harkness~J.\,H. et al. --- sleep disruption and oxidative stress in PV cells; Vyazovskiy~V.\,V. et al., 2009 --- firing rate and sleep; He~B.\,J., 2014 --- the scale-free spectrum of brain activity.
\item Work of A.~Grace's group on the MAM model; work of O.~Howes on dopamine synthesis before psychosis; Do~K.\,Q. et al., 2000, and the 2021 synthesis on glutathione in patients.
\item A review of the structure and function of PV cells; work on the stability of the cortex (Tsodyks, 1997; Ozeki, 2009); studies of the duration of symptoms and of the prodrome.
\end{itemize}
The bibliographic details of the second group are given from working notes and require verification before citation.

\textbf{Textbook knowledge} (tag \tU) is written from memory. Before citing such a fact it should be checked against a textbook, for example: Kandel~E. et al., Principles of Neural Science; Purves~D. et al., Neuroscience; Alberts~B. et al., Molecular Biology of the Cell.
}
